% Generated draft: requires mathematical and source review.
\documentclass[11pt]{article}
\usepackage[margin=25mm]{geometry}
\usepackage{fontspec}
\setmainfont{Latin Modern Roman}
\usepackage{amsmath,amssymb,mathtools,unicode-math}
\setmathfont{Latin Modern Math}
\usepackage{enumitem,needspace,xurl,hyperref}
\setlength{\parindent}{0pt}
\setlength{\parskip}{5pt}
\setlength{\emergencystretch}{3em}
\newcommand{\N}{\mathbb{N}}
\newcommand{\Z}{\mathbb{Z}}
\newcommand{\Q}{\mathbb{Q}}
\newcommand{\R}{\mathbb{R}}
\newcommand{\C}{\mathbb{C}}
\newcommand{\F}{\mathbb{F}}
\newcommand{\A}{\mathbb{A}}
\newcommand{\PP}{\mathbb{P}}
\newcommand{\E}{\mathbb{E}}
\newcommand{\eps}{\varepsilon}
\newcommand{\dd}{\,\mathrm{d}}
\begin{document}
\section*{UnsolvedMath ID: 701335}
\textbf{Source problem code:} catalogue entry
\subsection{Definitions and mathematical statement}
Work in the language of groups, with multiplication, inversion and identity element. Let
\[
G=\langle a,b\mid a^{-1}ba=b^{-1}\rangle.
\]
A group quasi-identity is a universally quantified implication
\[
\forall x_1\cdots\forall x_k\;\left(\bigwedge_{i=1}^{m}u_i(x_1,\ldots,x_k)=v_i(x_1,\ldots,x_k)\;\Longrightarrow\;u(x_1,\ldots,x_k)=v(x_1,\ldots,x_k)\right),
\]
where \(k,m\) are nonnegative integers and all displayed terms are group terms; the empty conjunction is interpreted as true. Define the quasivariety generated by \(G\) by
\[
\mathcal Q(G)=\{H\mid H\text{ is a group satisfying every group quasi-identity valid in }G\}.
\]
\textbf{Question (A. I. Budkin, Kourovka Notebook Problem 17.33).} Does there exist an integer \(n\geq 1\) and a set \(\Sigma\) of group quasi-identities, each using only variables from \(\{x_1,\ldots,x_n\}\), such that for every group \(H\),
\[
H\in\mathcal Q(G)\quad\Longleftrightarrow\quad H\models\Sigma?
\]
The set \(\Sigma\) may be infinite: the requested finiteness is a uniform bound on the number of variables, not a finite number of axioms. No finiteness or finite generation assumption is imposed on \(H\).

The supplied statement reports that the answer is known for all other groups with one defining relation, but supplies no supporting references or further scope qualifications for that assertion. The registry labels this problem open. Its background also records an explicitly unverified AI claim of a negative answer; that claim is not adopted here as a theorem or as a verified resolution.
\subsection{Short English statement}
For the group \(\langle a,b\mid a^{-1}ba=b^{-1}\rangle\), ask whether its generated quasivariety can be axiomatized by possibly infinitely many quasi-identities with a common finite bound on their numbers of variables. The supplied open label and the explicitly unverified claim of a negative answer require separate treatment in mathematical review.
\subsection{Sources}
\begin{itemize}
\item Local UnsolvedMath registry, permanent ID \(701335\), source code \texttt{catalogue entry}, title \emph{Kourovka Notebook Problem 17.33}; supplied upstream revision \texttt{d371ee917ac9647fbb06cd6fc529deec6f3dda82}.
\item A. I. Budkin, Problem 17.33, \emph{The Kourovka Notebook (Unsolved Problems in Group Theory)}, 17th issue (2010). The supplied record also identifies the 21st issue edition (2026), arXiv:1401.0300v48. The registry describes the collection as Issues 1--21 (1965--2026), edited by E. I. Khukhro and V. D. Mazurov. These are supplied bibliographic attributions, not independently checked citations.
\item Supplied background and status-review note: an entry marked \(\#\) in the 21st issue (October 2026) reports an explicitly unverified AI claim of a negative answer, attributed to M. Kinyon and J. Urban, \emph{Preprint} (2026), with the supplied link \url{https://kourovkanotebookorg.wordpress.com/wp-content/uploads/2026/09/kourovka-mathematics.pdf}. No verification of this claim or its attribution is asserted here.
\item Supplied status review dated \texttt{2026-10-16}: canonical status \texttt{open}, evidence category \texttt{unrefereed\_claim}. This is supplied dated curation, not an independently verified literature assessment.
\end{itemize}
\end{document}
